\section{Environments and Robots}
\label{sec:envs}
% Source: docs/ENVS.md and docs/ROBOTS.md on branch envs-r7.

\paragraph{Simulator and integrity.} All environments are procedural MuJoCo \citep{todorov2012mujoco} scenes
built from primitives and procedural meshes. Objects move only through robot actuator torques and forces (and
an environment-owned conveyor drive). An integrity monitor checks at every physics step that no object pose or
velocity was written, no external force was applied to an object, and no weld is active on an object;
violations are counted per episode, and the data-generation runs report none. Suction is a contact-mediated
adhesion actuator with capped force.

\paragraph{Observations.} The actor sees only a sensor stream with per-episode randomised noise, latency and
dropout: a fixed-camera pose tracker with line-of-sight occlusion (10\,Hz), joint encoders, a tared wrist
force/torque sensor, gripper state and tactile pads (25\,Hz), and one-shot size and shape estimates. From
round~3 the stream also carries two actor-legal fragility cues, a vision material reading and a tactile squeeze
stiffness reading. True hidden physics is available to the evaluator only.

\paragraph{Task families.} The full suite has 18 families across bulk handling, fine manipulation and one
deformable task. Each has a success predicate on true state and hidden variables chosen to change the best
action. Table~\ref{tab:families} lists the families used in this paper. Scenes are ``site-dressed'' from
round~3 on (warehouse hall, racking, lighting; only nearby props collide), with dressing drawn from a separate
random stream so that object poses and hidden physics are unchanged.

\begin{table}[!htbp]
  \centering\small
  \caption{Task families referred to in this paper (a subset of the suite).}
  \label{tab:families}
  \begin{tabular}{@{}l p{0.36\linewidth} p{0.30\linewidth}@{}}
    \toprule
    family & success & hidden variables that matter \\
    \midrule
    \fm{bin_clearing}   & every part in the tote (lot protocol: and undamaged) & mass, grip friction, crush limit \\
    \fm{conveyor_sort}  & every part in the bin of its class & belt speed, mass, grip friction \\
    \fm{depalletize}    & all cartons upright on the outfeed, no drops & carton mass (hidden heavy cartons) \\
    \fm{fixture_place}  & within 2\,mm and $3^\circ$ of the mark, released & mass, grip friction \\
    \fm{pouch_transfer} & deformable pouch in the tote, released & stiffness, mass, friction \\
    \fm{tool_return}    & tool shank seated in its holder & head mass and imbalance, clearance \\
    \fm{tight_insert}   & square peg (0.1--0.5\,mm clearance) seated past snap-fit lips & lip stiffness, dent limit \\
    \fm{fastener_pick}  & ordered screws and connector blocks moved from a tray into a kit & not characterised here \\
    \fm{drawer_open} (held out) & drawer opened to 85\% of travel & friction, damping, contents mass \\
    \bottomrule
  \end{tabular}
\end{table}

\paragraph{Robots.} Arms are driven by torque only, through a task-space inverse-dynamics (operational-space)
controller \citep{khatib1987unified} with torque-limit scaling, and a library of parameterised skills (reach,
pick, place, push, insert, open, probe). The main embodiments are a UR5e with a Robotiq 2F-85 gripper and a
Franka Panda; a Kinova Gen3 with a 2F-85 is the \emph{held-out robot}, absent from every training and tuning
set. A Cartesian gantry (``probe'' robot) with force limits and the same skill interface serves as a
high-competence reference embodiment; gantry results are not arm results and are reported separately.
Training data additionally includes other embodiments (suction arms, a bimanual platform, a mobile manipulator,
low-cost arms). A fixed-base Unitree G1 humanoid is used only in the sprint's push demonstration.

\paragraph{Splits.} Train, validation and test are split by root seed (80/10/10), decided before generation.
Held-out splits use test-bucket seeds: \fm{heldout_robot} (Gen3) and a held-out family (\fm{drawer_open},
excluded from training and validation data from round~3).

\paragraph{Lot protocol.} In the lot experiments on \fm{bin_clearing}, all parts in a bin come from one lot
that shares a density factor, grip friction and crush limit (fragile with probability 0.6). The actor is told
only that the parts are one lot. Fragile lots are drawn so that a policy knowing the lot can always succeed,
but the admissible grip window is often narrow and requires a slower lift.

\paragraph{Scripted baselines.}
\emph{Nominal}: a conservative engineer's heuristic with fixed parameters and escalating retries; it never
reads labels. \emph{Probe-then-act (\pta)}: a light 4\,N squeeze, then a grip of at most 0.7 times the estimated
crush limit. \emph{Force-ramp} (\fm{tight_insert} only): press at 4\,N, add 2\,N per stall, up to 14\,N.
\emph{Oracle-parameter} script: the same skills with parameters set from the true hidden physics (labelled).

\begin{figure}[!htbp]
  \centering
  \maybefig[0.5\linewidth]{../results/v2/envs/contact_sheet_probe.png}{contact sheet of the task families}
  \caption{Contact sheet of the task families (gantry embodiment, undressed scenes).}
  \label{fig:families}
\end{figure}
